\section{SmallWood Proof-Stack Details}
\label{app:smallwood}

This appendix records the compiled surfaces used by the issuance and showing statements.  It is descriptive: the
semantic statements are those of Section~\ref{sec:smallwood-model}.

\subsection{Row inventories}
\paragraph{Issuance rows.}
The issuance proof commits rows encoding:
\[
  \vecm,
  \veck,
  r_0^H,
  r_1^H,
  \bar r,
  r_0,
  r_1,
  Z,
\]
where the long target-randomizer satisfies $r_0\in\Rq^{\ell_{r_0}}$.  The proof also carries carry rows for the
centering equations and auxiliary rows for range membership.  The active relations are
\begin{align*}
  \vecc&=\mACom[\vecm\Vert\veck\Vert r_0^H\Vert r_1^H\Vert\bar r]^\top,\\
  r_0&=\centerfunc(r_0^H+r_0^I),\qquad r_1=\centerfunc(r_1^H+r_1^I),\\
  (B_3-r_1)\Had Z&=1,\\
  T&=B_0+B_1(\vecm\Vert\veck)+B_2r_0+Z.
\end{align*}

\paragraph{Showing rows.}
The showing proof commits rows encoding:
\[
  \vecu,
  \vecm,
  \veck,
  r_0,
  r_1,
  Z,
\]
with the same dimension convention $r_0\in\Rq^{\ell_{r_0}}$, plus rows for the signature shortness gadget and the
PRF companion relation.  The active algebraic relations are
\begin{align*}
  (B_3-r_1)\Had Z&=1,\\
  \mA\vecu&=B_0+B_1(\vecm\Vert\veck)+B_2r_0+Z,\\
  \prftag&=\PRF(\veck,\nonce).
\end{align*}

\subsection{Commitment and centering constraints}
The Ajtai opening constraint is linear over $\Rq$ and is compiled coefficient-wise into SmallWood rows.  The centering
map is represented by carry variables.  With $\Delta=2\BoundMsg+1$, each coefficient satisfies
\[
  r_b = r_b^H+r_b^I-\Delta j_b,
  \qquad j_b\in\{-1,0,1\},
\]
for $b\in\{0,1\}$.  Membership in the bounded interval and in $\{-1,0,1\}$ is enforced by the vanishing
polynomials of Appendix~\ref{app:defs}.

\subsection{BB-tran constraints}
The BB-tran constraints are expressed in witness form.  Issuance checks the public target equation
\[
  T=B_0+B_1(\vecm\Vert\veck)+B_2r_0+Z
\]
and the inverse equation
\[
  (B_3-r_1)\Had Z=1.
\]
Showing uses the same inverse equation and replaces the public target by the signature surface $\mA\vecu$:
\[
  \mA\vecu=B_0+B_1(\vecm\Vert\veck)+B_2r_0+Z.
\]
Equivalently, the showing proof implies the eliminated coefficient-domain relation
\[
  (B_3-r_1)\Had\bigl(\mA\vecu-B_0-B_1(\vecm\Vert\veck)-B_2r_0\bigr)=1.
\]

\subsection{PRF companion constraints}
The PRF companion relation binds the public pair $(\nonce,\prftag)$ to the hidden signed key $\veck$.  The proof
commits key-lane rows, selected S-box checkpoint rows, and any auxiliary rows needed to replay the Poseidon2-style
linear layers between checkpoints.  Selector constraints tie the key-lane rows to the signed key coordinates; final
feed-forward and truncation constraints tie the public tag to the resulting PRF output.

\subsection{Single committed witness}
All rows in one proof live under one SmallWood commitment.  This ensures that the prover cannot use one hidden
message for the Ajtai opening, a different message for the BB-tran target, and a third key for the PRF.  In the showing
proof, the same hidden key $\veck$ is used simultaneously in the signature relation and the tag relation.
